% concatenated from hsyB20260814.tex
%%% 2026.07.08
\documentclass[12pt]{amsart}
\usepackage{amsthm,amsmath,amssymb,enumerate}
\addtolength{\textwidth}{2cm}
\addtolength{\oddsidemargin}{-1cm}
\addtolength{\evensidemargin}{-1cm}
\addtolength{\textheight}{3cm}
\addtolength{\topmargin}{-1.5cm}

\newtheorem{thm}{Theorem}[section]
\newtheorem{prop}[thm]{Proposition}
\newtheorem{lemm}[thm]{Lemma}
\newtheorem{coro}[thm]{Corollary}
\theoremstyle{definition}
\newtheorem*{rema}{Remark}

\newcommand{\Z}{\mathbb{Z}}
\newcommand{\Aut}{\operatorname{Aut}}
\newcommand{\Hol}{\operatorname{Hol}}
\newcommand{\ci}{\circ}

\begin{document}
\title[Classification of skew left braces]
{Classification of skew left braces with additive group
  isomorphic to the infinite dihedral group}

\author[A. Hanaki]{Akihide Hanaki}
\address{Faculty of Science, Shinshu University,
  Matsumoto 3-1-1, 390-8621, Japan}
\email{hanaki@shinshu-u.ac.jp}

\author[Y. Sakata]{Yuto Sakata}
\address{1-6-35 Miyanomori, Muroran,
  Hokkaido, 050-0073, Japan}
\email{yuto.sakata.math@gmail.com}

\author[H. Yoshino]{Hiroki Yoshino}
\address{2642 Okitsu, Katsuura, Chiba 299-5245, Japan}
\email{yoshino.math@gmail.com}

\begin{abstract}
  We classify all skew left braces with additive group
  isomorphic to the infinite dihedral group.
  There are ten isomorphism classes.
\end{abstract}

\maketitle

\section{Introduction}
Let $A$ be a set with two binary operations
$(a,b)\mapsto a\cdot b$ and $(a,b)\mapsto a \ci b$.
The triple $(A,\cdot,\ci)$ is called a \emph{skew left brace} if
$(A,\cdot)$ and $(A,\ci)$ are groups and
$$a\ci(b\cdot c)=(a\ci b)\cdot a^{-1} \cdot (a\ci c)$$
holds for all $a,b,c\in A$,
where $a^{-1}$ is the inverse element of $a$
with respect to the operation of $(A,\cdot)$.
The groups $(A,\cdot)$ and $(A,\ci)$ are called the \emph{additive group}
and the \emph{multiplicative group} of the skew left brace, respectively.
We remark that the additive group is not necessarily an abelian group.
(If $(A,\cdot)$ is abelian, then $(A,\cdot,\ci)$ is called a \emph{left brace}.)
We often abbreviate $a\cdot b$ as $ab$.
The identity element of $(A,\cdot)$ is also the identity element of $(A,\ci)$.

Let $G$ be a group.
It is natural to ask how many skew left braces with additive (multiplicative) group $G$ exist.
In this direction, there are many results for finite groups $G$
\cite{MR4414157, MR4391827, MR4815260, MR4678589, MR4113853,
  MR4513787, MR4950651, MR3647970}.
For the infinite cyclic group $G$ as an additive group,
it is known that there are exactly two isomorphism
classes of skew left braces \cite[Proof of Proposition 6]{MR2298848}.
In this paper, we determine all isomorphism classes of skew left braces
with additive group isomorphic to the infinite dihedral group $D_\infty$.
There are ten isomorphism classes of such skew left braces
(Theorem \ref{cls_slb}).
This answers a problem raised by Vendramin
\cite[after Problem 2.27]{MR3974481}.

It is known that skew left braces with additive group $G$ correspond to
regular subgroups of the holomorph of $G$, and they correspond to
maps $\gamma:A\to \Aut(A,\cdot)$ with some property
\cite[Theorem 2.2]{MR4089566}, \cite[Theorem 4.2]{MR3647970}.
We will consider these maps to classify skew left braces.
Our method is very simple,
but it requires a lot of case analysis and calculations.

%%%%%%%%%%%%%%%%%%%%%%%%%%%%%%%%%%%%%%%
\section{Preliminaries}
\subsection{Holomorphs and $\gamma$-maps}
Let $G$ be a group, and consider the automorphism group $\Aut(G)$.
The semidirect product of $G$ and $\Aut(G)$ is called
the \emph{holomorph} of $G$, and denoted by $\Hol(G)$.
Namely,
$$\Hol(G)=G\rtimes \Aut(G)=\{(a,f)\mid a\in G,\ f\in \Aut(G)\},$$
and the product is defined by
$$(a,f)(b,g)=(af(b), fg).$$
The holomorph acts on $G$ from the left by
$(a,f)x=af(x)$ for $(a,f)\in \Hol(G)$ and $x\in G$.

Let $G$ be a group.
A map $\gamma:G\to\Aut(G)$ is called a \emph{$\gamma$-map} on $G$ if
$$\gamma(a\gamma(a)(b))=\gamma(a)\gamma(b)$$
holds for all $a, b\in G$.
The automorphism group $\Aut(G)$ acts on the set of all $\gamma$-maps by
$$\gamma^\varphi(a)=\varphi^{-1}\gamma(\varphi(a))\varphi\quad
(a\in G,\ \varphi\in\Aut(G)).$$
We define an equivalence relation $\sim$ by the above action.

Now we are ready to give a correspondence between skew left braces
with additive group $G$ and some other objects
\cite[Theorem 2.2]{MR4089566}, \cite[Theorem 4.2]{MR3647970}.

\begin{prop}\label{prop2.1}
  Let $G$ be a group.
  There are bijections between the following sets:
  \begin{enumerate}[{\rm (1)}]
    \item The set of isomorphism classes of skew left braces
    with additive group $G$.
    \item The set of conjugacy classes of regular subgroups of $\Hol(G)$.
    \item The $\sim$ equivalence classes of $\gamma$-maps on $G$. 
  \end{enumerate}
\end{prop}

For any $\gamma$-map $\gamma:G\to \Aut(G,\cdot)$,
the binary operation
$$a\circ b = a \cdot \gamma(a)(b) \qquad (a,b\in G)$$
makes $(G,\cdot,\circ)$ into a skew left brace.

\subsection{The infinite dihedral group $D_\infty$}
The infinite dihedral group is
$$D_\infty=\langle x,y\mid y^2=1,\ yx=x^{-1}y\rangle.$$
We denote $x^ay^s$ by $(a,s)$, and so
$$D_\infty=\{(a,s)\mid a\in \Z,\ s\in \Z_2\},$$
where $\Z_2=\Z/2\Z$, with operation
$$(a,\ s)(b,\ t)=(a+(-1)^s b,\ s+t).$$

For the automorphism group of $D_\infty$,
we have the following proposition.

\begin{prop}\label{prop2.2}
  We have
  $\Aut(D_\infty)=\{[b,t]\mid b\in \Z,\ t\in\Z_2\}$,
  where $$[b,t](a,s)=((-1)^ta+bs,s).$$
  The following equations hold.
  \begin{enumerate}[{\rm (1)}]
    \item $[a,s][b,t] = [a+(-1)^s b, s+t]$,
    \item $[a,s]^{-1}=[(-1)^{s+1}a,s]$,
    \item $[b,t]^{-1}[a,s][b,t]=[(-1)^ta+(-1)^{t+1}b+(-1)^{s+t}b,s]$.
  \end{enumerate}
\end{prop}

We remark that $s\in \Z_2$ can appear in
the first entry of $[b,t](a,s)=((-1)^ta+bs,s)$.
The elements of $\Z_2=\{0,1\}$ are naturally embedded into $\Z$. 
We will use the similar notation throughout this paper.

%%%%%%%%%%%%%%%%%%%%%%%%%%%%%%%%%%%%%%%
\section{Main results -- skew left braces with additive group $D_\infty$}
We classify all skew left braces with additive group $D_\infty$.
By Proposition \ref{prop2.1},
we classify equivalence classes of $\gamma$-maps on $D_\infty$.
We determine all $\gamma$-maps, and then we consider their equivalence.
Consequently we will get ten isomorphism classes of such skew left braces.
The following theorem is the main result of this paper.

\begin{thm}\label{cls_slb}
  Let $(A,\cdot,\ci)$ be a skew left brace with $(A,\cdot)=D_\infty$.
  Then $(A,\cdot,\ci)$ is isomorphic to a skew left brace
  defined by one of the following multiplications.
  Conversely, the following ten operations define
  non-isomorphic skew left braces.
  \begin{itemize}
    \item[(1)] $(a,s)\circ(b,t)=(a,s)(b,t)=(a+(-1)^sb,s+t)$
    (the trivial skew brace)
    \item[(2,0)] $(a,s)\circ(b,t)=(a+(-1)^sb-2at,s+t)$
    \item[(2,1)] $(a,s)\circ(b,t)=(a+(-1)^sb-2at+(-1)^s st,s+t)$
    \item[(3)] $(a,s)\circ(b,t)=(a+b-2at,s+t)$
    \item[(4,0)] $(a,s)\circ(b,t)=(a+b,s+t)$
    \item[(4,1)] $(a,s)\circ(b,t)=(a+b+(-1)^s st,s+t)$
    \item[(5,0)] $(a,s)\circ(b,t)=\begin{cases}
      (a+(-1)^sb,s+t), & \text{if $a$ is even},\\
      (a+(-1)^{s+1}b,s+t), & \text{if $a$ is odd}.
    \end{cases}$
    \item[(5,1)] $(a,s)\circ(b,t)=\begin{cases}
      (a+(-1)^sb,s+t), & \text{if $a$ is even},\\
      (a+(-1)^{s+1}b+(-1)^s2t,s+t), & \text{if $a$ is odd}.
    \end{cases}$
    \item[(6,0)] $(a,s)\circ(b,t)=\begin{cases}
      (a+b-2at,s+t), & \text{if $a+s$ is even},\\
      (a-b-2at+2t,s+t), & \text{if $a+s$ is odd}.
    \end{cases}$
    \item[(6,1)] $(a,s)\circ(b,t)=\begin{cases}
      (a-2at+(-1)^s b + 4st,s+t), & \text{if $a$ is even},\\
      (a-2(a-1)t+(-1)^{s+1}b+(-1)^s2t,s+t), & \text{if $a$ is odd}.
    \end{cases}$
  \end{itemize}
\end{thm}

Theorem \ref{cls_slb} will be proved by applying $a\ci b=a\gamma(a)(b)$ to
the classification of $\gamma$-maps (Theorem \ref{cls_gamma}).
For the structure of the multiplicative groups $(A,\ci)$,
we have the following corollary.

\begin{coro}
  The multiplicative groups $(A,\ci)$ of the skew left braces
  $(A,\cdot,\ci)$ in Theorem \ref{cls_slb} are isomorphic to 
  \begin{itemize}
    \item $\Z$ for (2,1), (4,1),
    \item $\Z\times \Z_2$ for (2,0), (4,0), 
    \item $D_\infty=\Z\rtimes \Z_2$ for (1), (3), (5,0), (6,0).
    \item $D_\infty\times \Z_2=(\Z\rtimes \Z_2)\times \Z_2$ for (5,1), (6,1).
  \end{itemize}
\end{coro}

\begin{proof}
  Write $\Z=\langle x\rangle$.
  For (2,1) or (4,1), $x\mapsto(0,1)$  gives an isomorphism.
  
  Write $\Z\times \Z_2=\langle x\rangle\times \langle y\rangle$.
  For (2,0) or (4,0), $x\mapsto(1,0)$, $y\mapsto(0,1)$
  gives an isomorphism.
  
  Write $D_\infty=\Z\rtimes \Z_2=\langle x\rangle\rtimes\langle y\rangle$.
  For (1) or (3), $x\mapsto(1,0)$, $y\mapsto(0,1)$ gives an isomorphism.
  For (5,0), $x\mapsto(1,1)$, $y\mapsto(0,1)$ gives an isomorphism.
  For (6,0), $x\mapsto(0,1)$, $y\mapsto(1,0)$ gives an isomorphism.

  Write $D_\infty\times \Z_2=(\Z\rtimes \Z_2)\times \Z_2
  =(\langle x\rangle\rtimes\langle y\rangle)\times\langle z\rangle$.
  For (5,1), $x\mapsto(2,0)$, $y\mapsto(1,0)$, $z\mapsto(1,1)$
  gives an isomorphism.
  For (6,1), $x\mapsto(0,1)$, $y\mapsto(1,0)$, $z\mapsto(2,1)$
  gives an isomorphism.
\end{proof}

%%%%%%%%%%%%%%%%%%%%%%%%%%%%%%%%%%%%%%%
\section{$\gamma$-Maps}
In this section, we will determine all $\gamma$-maps on $D_\infty$.
We will consider their equivalence in the next section.
For a $\gamma$-map $\gamma$, we write
$$\gamma(a,s)=[p(a,s),q(a,s)],$$
where $p:D_\infty\to \Z$ and $q:D_\infty\to \Z_2$.
Easily, we can see that $p(0,0)=0$ and $q(0,0)=0$.

\begin{prop}\label{prop_gamma}
  A map $\gamma:D_\infty\to \Aut(D_\infty)$ is a $\gamma$-map if and only if
  the following two equations hold for all $(a,s), (b,t)\in D_\infty$:
  \begin{align}
    p(a,s)+(-1)^{q(a,s)}p(b,t)
    &=p(a+(-1)^{s+q(a,s)}b+(-1)^sp(a,s)t,s+t),\tag{$\star$}\\
    q(a,s)+q(b,t)
    &=q(a+(-1)^{s+q(a,s)}b+(-1)^sp(a,s)t,s+t).\tag{$\star\star$}
  \end{align}
  We remark that the second equation is in $\Z_2$.
\end{prop}

\begin{proof}
  The result holds by the definition of $\gamma$-maps.
\end{proof}

We divide our arguments into two cases:
the case $q(1,0)=0$ and the case $q(1,0)=1$.

%%%%%%%%%%%%%%%%%%%%%%%%%%%%%%%%%%%%%%%
\subsection{The case $q(1,0)=0$}
In this subsection, we assume $q(1,0)=0$.

\begin{lemm}\label{lem4.2}
  $p(b,0)=bp(1,0)$, $q(b,0)=0$,
  $p(b,1)=p(0,1)-bp(1,0)$, and $q(b,1)=q(0,1)$
  for all $b\in \Z$.
\end{lemm}

\begin{proof}
  Set $(a,s)=(1,0)$ and $t=0$ for Proposition \ref{prop_gamma}.
  By $q(1,0)=0$, we have
  $p(1,0)+p(b,0)=p(1+b,0)$ and $q(1,0)+q(b,0)=q(1+b,0)$.
  Thus $p(b,0)=bp(1,0)$ and $q(b,0)=0$.

  Set $(a,s)=(0,1)$ and $t=0$ for Proposition \ref{prop_gamma} ($\star$).
  We have
  $p(0,1)+(-1)^{q(0,1)}bp(1,0)=p((-1)^{1+q(0,1)}b,1)$.
  Thus $p(b,1)=p(0,1)-bp(1,0)$.

  Set $(a,s)=(0,1)$ and $t=0$ for Proposition \ref{prop_gamma} ($\star\star$).
  We have
  $q(0,1)+q(b,0)=q((-1)^{1+q(0,1)}b,1)$.
  Thus $q(b,1)=q(0,1)+q((-1)^{1+q(0,1)}b,0)=q(0,1)$.
\end{proof}

\begin{lemm}\label{lem4.3}
  $p(1,0)=0$ or $-2$. Moreover,
  \begin{enumerate}[{\rm (1)}]
    \item if $q(0,1)=0$,
    then $(p(1,0),p(0,1))=(0,0)$ or $(-2, k)$ for some $k\in \Z$,
    \item if $q(0,1)=1$,
    then $(p(1,0),p(0,1))=(-2,0)$ or $(0, k)$ for some $k\in \Z$.
  \end{enumerate}
\end{lemm}

\begin{proof}
  Set $(a,s)=(1,0)$ and $(b,t)=(0,1)$ for ($\star$).
  We have $p(1,0)+p(0,1)=p(1+0+p(1,0),1)$ and so
  $p(1,0)(2+p(1,0))=0$. This shows $p(1,0)=0$ or $-2$.
  
  Set $(a,s)=(b,t)=(0,1)$ for ($\star$).
  We have $p(0,1)+(-1)^{q(0,1)}p(0,1)=p(-p(0,1),0)=-p(0,1)p(1,0)$.
  If $q(1,0)=0$, then $p(0,1)(2+p(1,0))=0$
  and this means $p(0,1)=0$ or $p(1,0)=-2$.
  If $q(1,0)=1$, then $p(0,1)p(1,0)=0$
  and this means $p(0,1)=0$ or $p(1,0)=0$.
\end{proof}

\begin{prop}\label{prop00}
  Suppose $\gamma$ is a $\gamma$-map and $q(1,0)=q(0,1)=0$.
  Then one of the following statements holds:
  \begin{enumerate}[{\rm (1)}]
    \item $p(a,s)=0$, $q(a,s)=0$,
    \item there exists $k\in\Z$ such that
    $p(a,s)=sk+(-1)^{s+1}2a$, $q(a,s)=0$.
  \end{enumerate}
  Moreover, they satisfy the conditions in Proposition \ref{prop_gamma}.
\end{prop}

\begin{proof}
  Consider the case in Lemma \ref{lem4.3} (1).
  By Lemma \ref{lem4.2}, $q(a,s)=0$.
  If $(p(1,0),p(0,1))=(0,0)$, then (1) holds.
  If $(p(1,0),p(0,1))=(-2,k)$, then
  we can write $p(a,s)=sk+(-1)^{s+1}2a$.

  For (1),  the conditions in Proposition \ref{prop_gamma} are clear.
  We consider ($\star$) for (2).
  \begin{align*}
    \mathrm{LHS} &= sk+(-1)^{s+1}2a+tk+(-1)^{t+1}2b,\\
    \mathrm{RHS} &= p(a+(-1)^sb+(-1)^sp(a,s)t,s+t)\\
                 &= p(a+(-1)^sb+(-1)^s(sk+(-1)^{s+1}2a)t,s+t)\\
                 &= (s+t)k+(-1)^{s+t+1}2(a+(-1)^sb+(-1)^s(sk+(-1)^{s+1}2a)t)\\
                 &= (s+t)k+(-1)^{s+t+1}2a+(-1)^{t+1} 2b
                   +(-1)^{t+1}2stk+(-1)^{s+t}4at.
  \end{align*}
  We remark that $(s+t)=0$ if $s=t=1$.
  
  If $t=0$, then
  $$\mathrm{LHS} = sk+(-1)^{s+1}2a-2b= \mathrm{RHS}.$$ 
  If $(s,t)=(0,1)$, then
  \begin{align*}
    \mathrm{LHS} &= -2a+k+2b,\\
    \mathrm{RHS} &= k+2a+2b-4a = -2a+k+2b.
  \end{align*}
  If $(s,t)=(1,1)$, then
  \begin{align*}
    \mathrm{LHS} &= 2k+2a+2b,\\
    \mathrm{RHS} &= -2a+2b+2k+4a=2k+2a+2b.
  \end{align*}
  For every case, the equation ($\star$) holds.
  The equation ($\star\star$) for (2) is clear.
\end{proof}

\begin{prop}\label{prop01}
  Suppose $\gamma$ is a $\gamma$-map and $q(1,0)=0$, $q(0,1)=1$.
  Then one of the following statements holds:
  \begin{enumerate}[{\rm (1)}]
    \item $p(a,s)=(-1)^{s+1}2a$, $q(a,s)=s$,
    \item $p(a,s)=sk$, $q(a,s)=s$.
  \end{enumerate}
  They satisfy the conditions in Proposition \ref{prop_gamma}.
\end{prop}
  
\begin{proof}
  Consider the case in Lemmas \ref{lem4.2} and \ref{lem4.3} (2).
  We obtain (1) and (2).

  Consider the condition ($\star$) for (1).
  \begin{align*}
    \mathrm{LHS} &= (-1)^{s+1}2a+(-1)^s(-1)^{t+1}2b\\
                 &= (-1)^{s+1}2a+(-1)^{s+t+1}2b,\\
    \mathrm{RHS} &= p(a+b-2at,s+t)\\
                 &= (-1)^{s+t+1}2(a+b-2at)\\
                 &= (-1)^{s+t+1}2a+(-1)^{s+t+1}2b+(-1)^{s+t}4at\\
                 &= (-1)^{s+t}2a(-1+2t)+(-1)^{s+t+1}2b.
  \end{align*}
  Thus ($\star$) holds, and ($\star\star$) for (1) is clear.

  For the equation ($\star$) for (2),
  \begin{align*}
    \mathrm{LHS} &= sk+(-1)^stk,\\
    \mathrm{RHS} &= p(a+b+(-1)^skst,s+t)
                     = (s+t)k.
  \end{align*}
  We remark that $s+t=0$ if $s=t=1$.
  It is easy to verify that $\mathrm{LHS}= \mathrm{RHS}$.
  For ($\star\star$), clearly $\mathrm{LHS}=s+t=\mathrm{RHS}$ holds.
\end{proof}

%%%%%%%%%%%%%%%%%%%%%%%%%%%%%%%%%%%%%%% 
\subsection{The case $q(1,0)=1$}
In this subsection, we assume $q(1,0)=1$.

\begin{lemm}\label{lem4.6}
  $p(2\ell,0)=\ell p(2,0)$,
  $p(2\ell +1, 0)=\ell p(2,0)+p(1,0)$, and
  \begin{enumerate}[{\rm (1)}]
    \item if $q(0,1)=0$, then $q(a,s)=a$,
    \item if $q(0,1)=1$, then $q(a,s)=a+s$.
  \end{enumerate}
\end{lemm}

\begin{proof}
  Set $(a,s)=(-1,0)$ and $(b,t)=(1,0)$ for  ($\star\star$).
  We get $q(-1,0)+q(1,0)=q(-1+(-1)^{q(-1,0)},0)$.
  If $q(-1,0)=0$, then $0+1=q(-1+1,0)=q(0,0)=0$ and this is a contradiction.
  Thus $q(-1,0)=1$.
  
  Set $s=0$ and $(b,t)=((-1)^{q(a,0)},0)$  for ($\star\star$). 
  We have $q(a,0)+q((-1)^{q(a,0)},0)=q(a+1,0)$.
  Since $q(a,0)\in \{0, 1\}$ and $q(0,0)=0$,
  we can conclude that $q(a,0)=a$ (in $\Z_2$).

  Set $(a,s)=(2,0)$ and $t=0$ for ($\star$).
  We get $p(2,0)+p(b,0)=p(2+b,0)$.
  This shows $p(2\ell,0)=\ell p(2,0)$ and
  $p(2\ell +1, 0)=\ell p(2,0)+p(1,0)$.

  Set $(a,s)=(0,1)$ and $t=0$ for ($\star\star$). .
  We have $q(0,1)+q(b,0)=q((-1)^{1+q(0,1)}b,1)$.
  If $q(0,1)=0$, then $b=q(b,0)=q(-b,1)$ and this shows $q(a,s)=a$.
  If $q(0,1)=1$, then $1+b=1+q(b,0)=q(b,1)$ and this shows $q(a,s)=a+s$.
\end{proof}

\begin{lemm}\label{lem4.7}
  $p(a,s)$ is even for any $(a,s)\in D_\infty$.
\end{lemm}

\begin{proof}
  Set $(b,t)=(0,1)$ for ($\star\star$).
  We have $q(a,s)+q(0,1)=q(a+(-1)^sp(a,s),1+s)$ in $\Z_2$.
  If $q(0,1)=0$, then $a=q(a,s)=q(a+(-1)^sp(a,s),1+s)=a+(-1)^sp(a,s)$
  and this shows $p(a,s)$ is even.
  If $q(0,1)=1$, then $a+s+1=q(a,s)+1=q(a+(-1)^sp(a,s),1+s)=a+s+1+(-1)^sp(a,s)$
  and this shows $p(a,s)$ is even.
\end{proof}

We consider the case $q(0,1)=0$.

\begin{lemm}\label{lem4.8}
  If $q(0,1)=0$, then $p(2\ell,1) = p(0,1)-\ell p(2,0)$,
  $p(2\ell-1,1)= p(0,1)-\ell p(2,0)+p(1,0)$,
  and $p(2,0)=0$ or $-4$.
\end{lemm}

\begin{proof}
  Set $(a,s)=(0,1)$ and $t=0$ for ($\star$).
  We have $p(0,1)+p(b,0)=p(-b,1)$.
  With Lemma \ref{lem4.6}, we have $p(2\ell,1) = p(0,1)-\ell p(2,0)$, and
  $p(2\ell-1,1)= p(0,1)-\ell p(2,0)+p(1,0)$.

  Set $(a,s)=(2,0)$ and $(b,t)=(0,1)$ for ($\star$).
  We have $p(2,0)+p(0,1)=p(2+p(2,0),1)$ and so
  $p(2,0)(p(2,0)+4)=0$.
\end{proof}

\begin{prop}\label{prop10}
  If $q(1,0)=1$ and $q(0,1)=0$, then one of the following statements holds:
  \begin{enumerate}[{\rm (1)}]
    \item 
    $p(2\ell,s)= 0$,
    $p(2\ell-1,s) = 2k$, and
    $q(a,s)=a$ 
    for some $k\in \Z$,
    \item 
    $p(2\ell,0)= -4\ell$,
    $p(2\ell+1,0) = -4\ell+2k$,
    $p(2\ell,1) = 4\ell-2k-2$,
    $p(2\ell-1,1) =4\ell-2$, and
    $q(a,s)=a$
    for some $k\in \Z$.
  \end{enumerate}
  They satisfy the conditions in Proposition \ref{prop_gamma}.
\end{prop}

\begin{proof}
  Assume $p(2,0)=0$.
  Set $(a,s)=(0,1)$ and $(b,t)=(0,1)$ for ($\star$).
  By Lemmas, \ref{lem4.6}, \ref{lem4.7}, and \ref{lem4.8},
  $p(0,1)+p(0,1)=p(-p(0,1),0)=0$.
  Thus $p(0,1)=0$ and $q(a,s)=a$ hold for all
  $(a,s)\in D_\infty$.
  Since $p(1,0)$ is even, we can set $2k=p(1,0)$ for some $k\in\Z$. 
  We can determine all $p(a,s)$ as in (1)
  by Lemmas \ref{lem4.6} and \ref{lem4.8}.
  The condition ($\star\star$) holds by
  \begin{align*}
    \text{LHS} &= q(a,s)+q(b,t)=a+b,\\
    \text{RHS} &= q(a+(-1)^{s+q(a,s)}b+(-1)^sp(a,s)t,s+t)=a+b.
  \end{align*}
  We can check the equation ($\star$) for all $(a,s),(b,t)\in D_\infty$.
  See Appendix A.

  Assume $p(2,0)=-4$.
  Set $(a,s)=(1,0)$ and $(b,t)=(0,1)$ for ($\star$).
  We have $p(1,0)-p(0,1)=p(0,1)+3p(1,0)+4$,
  and this means $p(0,1)=-p(1,0)-2$.
  Since $p(1,0)$ is even, we can set $2k=p(1,0)$ for some $k\in\Z$.
  The equations in (2) hold.
  The condition ($\star\star$) holds similarly to the above case.
  We can check the condition ($\star$) for all $(a,s),(b,t)\in D_\infty$.
  See Appendix B.
\end{proof}

We consider the case $q(0,1)=1$.

\begin{lemm}\label{lem4.10}
  If $q(0,1)=1$, then
  $p(2\ell,1)= p(0,1)-\ell p(2,0)$,
  $p(2\ell+1,1)= p(0,1)-\ell p(2,0)-p(1,0)$,
  and $p(2,0)=0$ or $-4$.
\end{lemm}

\begin{proof}
  Set $(a,s)=(0,1)$ and $t=0$ for ($\star$).
  We have $p(0,1)-p(b,0)=p(b,1)$.
  With Lemma \ref{lem4.6}, we can conclude
  $p(2\ell,1)= p(0,1)-\ell p(2,0)$ and
  $p(2\ell+1,1)= p(0,1)-\ell p(2,0)-p(1,0)$.
  
  Set $(a,s)=(2,0)$ and $(b,t)=(0,1)$ for ($\star$).
  By Lemma \ref{lem4.7},
  we have $p(2,0)+p(0,1)=p(2+p(2,0),1)=p(0,1)-(2+p(2,0))p(2,0)/2$,
  and this shows $p(2,0)=0$ or $-4$.
\end{proof}

\begin{prop}\label{prop11}
  If $q(1,0)=q(0,1)=1$, then one of the following statements holds:
  \begin{enumerate}[{\rm (1)}]
    \item  $p(2\ell,0) = 0$,
    $p(2\ell+1,0) = 2k$,
    $p(2\ell,1) = 2k$,
    $p(2\ell+1,1) = 0$, and
    $q(a,s)= a+s$ for some $k\in \Z$,
    \item
    $p(2\ell,0) = -4\ell$,
    $p(2\ell+1,0) = -4\ell+2k$,
    $p(2\ell,1) = 4\ell$,
    $p(2\ell+1,1) = 4\ell-2k$, and
    $q(a,s)= a+s$ for some $k\in \Z$.
  \end{enumerate}
\end{prop}

\begin{proof}
  Assume $p(2,0)=0$.
  Set $(a,s)=(1,0)$ and $(b,t)=(0,1)$ for ($\star$).
  We have $p(1,0)-p(0,1)=p(1+p(1,0),1)=p(0,1)-p(1,0)$.
  This shows that $p(1,0)=p(0,1)$.
  Since $p(1,0)$ is an even number, we set $p(1,0)=2k$.

  Remark that $p(a,s)=0$ if $a+s$ is even and $p(a,s)=2k$ if $a+s$ is odd.
  The left-hand side of ($\star$) is $0$ if $(a+s)+(b+t)$ is even and
  $1$ if $(a+s)+(b+t)$ is odd.
%  Similarly, we can see that the same holds for the right-hand side.
  Similarly, the same holds for the right-hand side.
  Similarly, the condition ($\star\star$) holds. 

  Assume $p(2,0)=-4$.
  Set $(a,s)=(1,0)$ and $(b,t)=(1,1)$ for ($\star$).
  We have $p(1,0)-p(1,1)=p(p(1,0),1)$ and this shows $p(0,1)=0$.
  Since $p(1,0)$ is even, we set $p(1,0)=2k$.

  For ($\star\star$),
  \begin{align*}
    \text{LHS} &= q(a,s)+q(b,t)=a+s+b+t,\\
    \text{RHS} &= q(a+(-1)^ab+(-1)^sp(a,s)t,s+t)=a+b+s+t,
  \end{align*}
  in $\Z_2$, and the equation holds.
  For the condition ($\star$), see Appendix C.
\end{proof}

%%%%%%%%%%%%%%%%%%%%%%%%%%%%%%%%%%%%%%% 
\subsection{All $\gamma$-maps}
Combining Propositions \ref{prop00}, \ref{prop01},
\ref{prop10} and \ref{prop11},
we list all $\gamma$-maps on $D_\infty$.

\begin{thm}
  $\gamma$-maps on $D_\infty$ are
  (in the following, $\ell$ is always an integer) :
  \begin{enumerate}[{\rm (1)}]
    \item
    $\gamma(a,s)=[0,0]$,
    \item
    $\gamma(a,s)=[sk+(-1)^{s+1}2a,0]$ for some $k\in \Z$,
    \item
    $\gamma(a,s)=[(-1)^{s+1}2a,s]$,
    \item
    $\gamma(a,s)=[sk,s]$ for some $k\in \Z$,
    \item 
    $\gamma(2\ell,s)=[0,0]$,
    $\gamma(2\ell+1,s)=[2k,1]$
    for some $k\in \Z$,
    \item
    $\gamma(2\ell,0)=[-4\ell,0]$,
    $\gamma(2\ell+1,0)=[-4\ell+2k,1]$,
    $\gamma(2\ell,1)=[4\ell-2k-2,0]$,
    $\gamma(2\ell-1,1)=[4\ell-2,1]$
    for some $k\in \Z$,
    \item
    $\gamma(2\ell,0)=[0,0]$,
    $\gamma(2\ell+1,0)=[2k,1]$,
    $\gamma(2\ell,1)=[2k,1]$,
    $\gamma(2\ell-1,1)=[0,0]$
    for some $k\in \Z$.
    \item
    $\gamma(2\ell,0)=[-4\ell,0]$,
    $\gamma(2\ell+1,0)=[-4\ell+2k,1]$,
    $\gamma(2\ell,1)=[4\ell,1]$,
    $\gamma(2\ell+1,1)=[4\ell-2k,0]$
    for some $k\in \Z$.
  \end{enumerate}
\end{thm}

For each case, we write $\gamma_{(i)}$ for the $\gamma$-map.
We write  $\gamma_{(i,k)}$ ($i\in\{2, 4, 5, 6, 7, 8\}$)
to specify the parameter $k$.

%%%%%%%%%%%%%%%%%%%%%%%%%%%%%%%%%%%%%%%
\section{Equivalence of $\gamma$-maps}
We determine the equivalences of $\gamma$-maps.

\begin{lemm}\label{lemCindy}
  $\gamma$-maps are characterized by the values at
  $(1,0)$, $(0,1)$, and $(2,0)$.
\end{lemm}

\begin{proof}
  We make a list for the values of $\gamma_{(i)}$. 
  $$
    \begin{array}{c||c|c|c}
      & \gamma_{(i)}(1,0) & \gamma_{(i)}(0,1) &\gamma_{(i)}(2,0)\\
      \hline
      \hline
      (1) & [0,0] & [0,0] & [0,0] \\
      \hline
      (2,k) & [-2,0] & [k,0] & [-4,0]\\
      \hline
      (3) & [-2,0] & [0,1] & [-4,0]\\
      \hline
      (4,k) & [0,0] & [k,1] & [0,0]\\
      \hline
      (5,k) & [2k,1] & [0,0] & [0,0]\\
      \hline
      (6,k) & [2k,1] & [-2k-2,0] & [-4,0]\\
      \hline
      (7,k) & [2k,1] & [2k,1] & [0,0]\\
      \hline
      (8,k) & [2k,1] & [0,1] & [-4,0]\\
      \hline
    \end{array}
  $$
  This shows the lemma.
\end{proof}

\begin{rema}
  Cindy Tsang pointed out that Lemma \ref{lemCindy}
  holds without the classification of $\gamma$-maps.
\end{rema}

\begin{thm}\label{cls_gamma}
  Representatives of $\gamma$-maps on $D_\infty$ are
  $$\gamma_{(1)}, \gamma_{(2,0)}, \gamma_{(2,1)}, \gamma_{(3)},
  \gamma_{(4,0)}, \gamma_{(4,1)}, \gamma_{(5,0)}, \gamma_{(5,1)},
  \gamma_{(6,0)}, \gamma_{(6,1)}. $$
\end{thm}

\begin{proof}
  We will show that
  \begin{itemize}
    \item $\{\gamma_{(1)}\}$ and $\{\gamma_{(3)}\}$
    are equivalence classes,
    \item for $i, j\in\{2,4,5,6\}$, $i\ne j$, $\gamma_{(i,k)}$
    and $\gamma_{(j,k')}$ are not equivalent,
    \item for $i\in \{2,4,5,6, 7, 8\}$ and $k, k'\in \Z$,
    $\gamma_{(i,k)}$ and $\gamma_{(i,k')}$ are equivalent
    if and only if $k\equiv k'\pmod{2}$, 
    \item $\gamma_{(5,0)}$ and $\gamma_{(7,1)}$ are equivalent, 
    \item $\gamma_{(5,1)}$ and $\gamma_{(7,0)}$ are equivalent,
    \item $\gamma_{(6,0)}$ and $\gamma_{(8,1)}$ are equivalent, and
    \item $\gamma_{(6,1)}$ and $\gamma_{(8,0)}$ are equivalent.
  \end{itemize}
  These prove the proposition.
  To show them, we will give a table of
  $\gamma_{(i)}^{[b,t]}(a,s)$ for $(a,s)=(1,0), (0,1), (2,0)$.

  {\small
  $$\begin{array}{c||c|c|c||c}
    & \gamma_{(i)}^{[b,t]}(1,0)
    & \gamma_{(i)}^{[b,t]}(0,1)
    & \gamma_{(i)}^{[b,t]}(2,0) & \gamma_{(j)}\\
    \hline
    \hline
    (1) & [0,0] & [0,0] & [0,0] & \gamma_{(1)}\\
    \hline
    (2,k) & [-2,0] & [(-1)^t(k+2b),0] & [-4,0]& \gamma_{(2,(-1)^t(k+2b))}\\
    \hline
    (3) & [-2,0] & [0,1] & [-4,0] & \gamma_{(3)}\\
    \hline
    (4,k) & [0,0] & [(-1)^t(k-2b),1] & [0,0] & \gamma_{(4,(-1)^t(k-2b))}\\
    \hline
    (5,k), \text{$b$:even} & [(-1)^t2(k-b),1] & [0,0] & [0,0] & \gamma_{(5,(-1)^t(k-b))}\\
    (5,k), \text{$b$:odd} & [(-1)^t2(k-b),1]  & [(-1)^t2(k-b),1] & [0,0]
                                & \gamma_{(7,(-1)^t(k-b))}\\
    \hline
    (6,k), \text{$b$:even, $t=0$} & [2(k-b),1] & [-2(k-b)-2,0] & [-4,0] &\gamma_{(6,k-b)}\\
    (6,k), \text{$b$:even, $t=1$} & [-2(k-b+2),1] & [2(k-b+2)-2,0] & [-4,0]&\gamma_{(6,-k+b-2)}\\
    (6,k), \text{$b$:odd, $t=0$} & [2(k-b),1] & [0,1] & [-4,0]&\gamma_{(8,k-b)}\\
    (6,k), \text{$b$:odd, $t=1$} & [-2(k-b+2),1] & [0,1] & [-4,0]&\gamma_{(8,-k+b-2)}\\
    \hline
    (7,k), \text{$b$:even}  & [(-1)^t2(k-b),1] & [(-1)^t2(k-b),1] & [0,0]
                                &\gamma_{(7,(-1)^t(k-b))}\\
    (7,k), \text{$b$:odd}  & [(-1)^t2(k-b),1] & [0,0] & [0,0]
                                &\gamma_{(5,(-1)^t(k-b))}\\
    \hline
    (8,k), \text{$b$:even, $t=0$}& [2(k-b),1] & [0,1] & [-4,0]&\gamma_{(8,k-b)}\\
    (8,k), \text{$b$:even, $t=1$}& [-2(k-b+2),1] & [0,1] & [-4,0]&\gamma_{(8,-k+b-2)}\\
    (8,k), \text{$b$:odd, $t=0$}& [2(k-b),1] & [-2(k-b)-2,0] & [-4,0]&\gamma_{(6,k-b)}\\
    (8,k), \text{$b$:odd, $t=1$}& [-2(k-b+2),1] & [2(k-b+2)-2,0] & [-4,0]&\gamma_{(6,-k+b-2)}\\
    \hline
  \end{array}$$}
  By the table, we can check all required conditions.
\end{proof}

Applying $a\ci b=a\gamma(a)(b)$ to Theorem \ref{cls_gamma}, we can show
our main result Theorem \ref{cls_slb}.

%%%%%%%%%%%%%%%%%%%%%%%%%%%%%%%%%%%%%%%
\section*{acknowledgment}
The authors would like to thank Cindy Tsang
and Yuta Kozakai for their valuable comments.
The first author was supported by JSPS KAKENHI Grant Number JP22K03266.

%%%%%%%%%%%%%%%%%%%%%%%%%%%%%%%%%%%%%%%
\bibliographystyle{amsplain}
\providecommand{\bysame}{\leavevmode\hbox to3em{\hrulefill}\thinspace}
\providecommand{\MR}{\relax\ifhmode\unskip\space\fi MR }
% \MRhref is called by the amsart/book/proc definition of \MR.
\providecommand{\MRhref}[2]{%
  \href{http://www.ams.org/mathscinet-getitem?mr=#1}{#2}
}
\providecommand{\href}[2]{#2}
\begin{thebibliography}{10}

\bibitem{MR4414157}
E.~Acri and M.~Bonatto, \emph{Skew braces of size {$p^2q$} {I}: {A}belian
  type}, Algebra Colloq. \textbf{29} (2022), no.~2, 297--320. %\MR{4414157}

\bibitem{MR4391827}
\bysame, \emph{Skew braces of size {$p^2q$} {II}: {N}on-abelian type}, J.
  Algebra Appl. \textbf{21} (2022), no.~3, Paper No. 2250062, 61. %\MR{4391827}

\bibitem{MR4815260}
A.~Ballester-Bolinches, R.~Esteban-Romero, and V.~P\'erez-Calabuig,
  \emph{Enumeration of left braces with additive group {$C_2\times C_2\times
  C_4\times C_4$}}, Ric. Mat. \textbf{73} (2024), no.~5, 2679--2685.
  %\MR{4815260}

\bibitem{MR4678589}
\bysame, \emph{Enumeration of left braces with additive group {$C_4\times
  C_4\times C_4$}}, Math. Comp. \textbf{93} (2024), no.~346, 911--919.
  %\MR{4678589}

\bibitem{MR4113853}
V.~G. Bardakov, M.~V. Neshchadim, and M.~K. Yadav, \emph{Computing skew left
  braces of small orders}, Internat. J. Algebra Comput. \textbf{30} (2020),
  no.~4, 839--851. %\MR{4113853}

\bibitem{MR4089566}
E.~Campedel, A.~Caranti, and I.~Del~Corso, \emph{Hopf-{G}alois structures on
  extensions of degree {$p^2q$} and skew braces of order {$p^2 q$}: the cyclic
  {S}ylow {$p$}-subgroup case}, J. Algebra \textbf{556} (2020), 1165--1210.
  %\MR{4089566}

\bibitem{MR4513787}
T.~Crespo, D.~Gil-Mu\~noz, A.~Rio, and M.~Vela, \emph{Left braces of size
  {$8p$}}, J. Algebra \textbf{617} (2023), 317--339. %\MR{4513787}

\bibitem{MR4950651}
\bysame, \emph{Determining skew left braces of size {$np$}}, Comm. Algebra
  \textbf{53} (2025), no.~11, 4817--4843. %\MR{4950651}

\bibitem{MR3647970}
L.~Guarnieri and L.~Vendramin, \emph{Skew braces and the {Y}ang-{B}axter
  equation}, Math. Comp. \textbf{86} (2017), no.~307, 2519--2534. %\MR{3647970}

\bibitem{MR2298848}
W.~Rump, \emph{Classification of cyclic braces}, J. Pure Appl. Algebra
  \textbf{209} (2007), no.~3, 671--685. %\MR{2298848}

\bibitem{MR3974481}
L.~Vendramin, \emph{Problems on skew left braces}, Adv. Group Theory Appl.
  \textbf{7} (2019), 15--37. %\MR{3974481}

\end{thebibliography}

%%%%%%%%%%%%%%%%%%%%%%%%%%%%%%%%%
\section*{Appendix A}
This is a suppliment for Proposition \ref{prop10} (1).
\medskip

{\tiny
\noindent\underline{$(a,s)=(2\ell,s)$, $(b,t)=(2m,t)$}
\begin{align*}
  \text{LHS} &= p(2\ell,s)+(-1)^{q(2\ell,s)}p(2m,t)=0,\\
  \text{RHS} &= p(2\ell+(-1)^{s+q(2\ell,s)}2m+(-1)^sp(2\ell,s)t,s+t)=0.
\end{align*}

\noindent\underline{$(a,s)=(2\ell,s)$, $(b,t)=(2m+1,t)$}
\begin{align*}
  \text{LHS} &= p(2\ell,s)+(-1)^{2\ell}p(2m+1,t)=2k,\\
  \text{RHS} &= p(2\ell+(-1)^{s+2\ell}(2m+1)+(-1)^sp(2\ell,s)t,s+t)=2k.
\end{align*}

\noindent\underline{$(a,s)=(2\ell+1,s)$, $(b,t)=(2m,t)$}
\begin{align*}
  \text{LHS} &= p(2\ell+1,s)+(-1)^{2\ell+1}p(2m,t)=2k\\
  \text{RHS} &= p(2\ell+(-1)^{s+2\ell+1}2m+(-1)^sp(2\ell+1,s)t,s+t)
                                   =2k.
\end{align*}

\noindent\underline{$(a,s)=(2\ell+1,s)$, $(b,t)=(2m+1,t)$}
\begin{align*}
  \text{LHS} &= p(2\ell+1,s)+(-1)^{2\ell+1}p(2m+1,t)
                                   =2k-2k=0,\\
  \text{RHS} &= p(2\ell+1+(-1)^{s+2\ell+1}(2m+1)+(-1)^sp(2\ell+1,s)t,s+t)
                                   =0.
\end{align*}
}

%%%%%%%%%%%%%%%%%%%%%%%%%%%%%%%%% 
\section*{Appendix B}
This is a suppliment for Proposition \ref{prop10} (2).
\medskip

{\tiny
\noindent\underline{$(a,s)=(2\ell,0)$, $(b,t)=(2m,0)$}
\begin{align*}
  \text{LHS} &= p(2\ell,0)+(-1)^{q(2\ell,0)}p(2m,0)=-4\ell-4m,\\
  \text{RHS} &= p(2\ell+2m,0)=-4\ell-4m.
\end{align*}

\noindent\underline{$(a,s)=(2\ell,0)$, $(b,t)=(2m+1,0)$}
\begin{align*}
  \text{LHS} &= p(2\ell,0)+p(2m+1,0)=-4\ell-4m+2k,\\
  \text{RHS} &= p(2\ell+(2m+1),0)=-4(\ell+m)+2k.
\end{align*}

\noindent\underline{$(a,s)=(2\ell,0)$, $(b,t)=(2m,1)$}
\begin{align*}
  \text{LHS} &= p(2\ell,0)+p(2m+1)=-4\ell+4m+2, \\
  \text{RHS} &= p(2\ell+2m+1+p(2\ell,0),1)=p(-2\ell+2m+1,1)= -4\ell+4m+2.
\end{align*}

\noindent\underline{$(a,s)=(2\ell,0)$, $(b,t)=(2m+1,1)$}
\begin{align*}
  \text{LHS} &= p(2\ell,0)+p(2m+1,1)=-4\ell+4m+2,\\
  \text{RHS} &= p(2\ell+(2m+1)-4\ell,1)=p(-2\ell+2m+1,1)= -4\ell+4m+2.
\end{align*}

\noindent\underline{$(a,s)=(2\ell+1,0)$, $(b,t)=(2m,0)$}
\begin{align*}
  \text{LHS} &= p(2\ell+1,0)-p(2m,0)=-4\ell+2k+4m,\\
  \text{RHS} &= p(2\ell+1-2m,0)==-4(\ell-m)+2k.
\end{align*}

\noindent\underline{$(a,s)=(2\ell+1,0)$, $(b,t)=(2m+1,0)$}
\begin{align*}
  \text{LHS} &= p(2\ell+1,0)-p(2m+1,0)=-4\ell+2k-(-4m+2k)=-4\ell+4m,\\
  \text{RHS} &= p(2\ell+1-(2m+1),0)==-4\ell+4m.
\end{align*}

\noindent\underline{$(a,s)=(2\ell+1,0)$, $(b,t)=(2m,1)$}
\begin{align*}
  \text{LHS} &= p(2\ell+1,0)-p(2m,1)=-4\ell+2k-(4m-2k-2)=-4\ell-4m+4k+2,\\
  \text{RHS} &= p(2\ell+1-2m+p(2\ell+1,0),1)= p(2\ell-2m+1-4\ell+2k,1)
                                   =-4\ell-4m+4k+2.
\end{align*}

\noindent\underline{$(a,s)=(2\ell+1,0)$, $(b,t)=(2m+1,1)$}
\begin{align*}
  \text{LHS} &= p(2\ell+1,0)-p(2m+1,1)=-4\ell+2k-(4m+2)=-4\ell-4m+2k-2,\\
  \text{RHS} &= p(2\ell+1-(2m+1)+(-4\ell+2k),1)= p(-2\ell-2m+2k,1)
                               =-4\ell-4m+2k-2.
\end{align*}

\noindent\underline{$(a,s)=(2\ell,1)$, $(b,t)=(2m,0)$}
\begin{align*}
  \text{LHS} &= p(2\ell,1)+p(2m,0)=4\ell-2k-2-4m,\\
  \text{RHS} &= p(2\ell-2m,1)=4\ell-4m-2k-2.
\end{align*}

\noindent\underline{$(a,s)=(2\ell,1)$, $(b,t)=(2m+1,0)$}
\begin{align*}
  \text{LHS} &= p(2\ell,1)+p(2m+1,0)=4\ell-4m-2\\
  \text{RHS} &= p(2\ell-(2m+1),1)=4\ell-4m-2.
\end{align*}

\noindent\underline{$(a,s)=(2\ell,1)$, $(b,t)=(2m,1)$}
\begin{align*}
  \text{LHS} &= p(2\ell,1)+p(2m,1)=4\ell-2k-2+4m-2k-2=4\ell+4m-4k-4,\\
  \text{RHS} &= p(2\ell-2m-(4\ell-2k-2),0)=p(-2\ell-2m+2k+2,0)
                                   =4\ell+4m-4k-4.
\end{align*}

\noindent\underline{$(a,s)=(2\ell,1)$, $(b,t)=(2m+1,1)$}
\begin{align*}
  \text{LHS} &= p(2\ell,1)+p(2m+1,1)=4\ell-2k-2+4m+2,\\
  \text{RHS} &= p(2\ell-(2m+1)-(4\ell-2k-2),0)
  = p(-2\ell-2m+2k+1,0)=4\ell+4m-4k+2k.
\end{align*}

\noindent\underline{$(a,s)=(2\ell+1,1)$, $(b,t)=(2m,0)$}
\begin{align*}
  \text{LHS} &= p(2\ell+1,1)-p(2m,0)=4\ell+2+4m,\\
  \text{RHS} &= p(2\ell-1+2m,1)=4\ell+4m+2.
\end{align*}

\noindent\underline{$(a,s)=(2\ell+1,1)$, $(b,t)=(2m+1,0)$}
\begin{align*}
  \text{LHS} &= p(2\ell+1,1)-p(2m+1,0)
                                   =4\ell+2-(-4m+2k)=4\ell+4m-2k+2,\\
  \text{RHS} &= p(2\ell+1+(2m+1),1)
  =4\ell+4m+4-2k-2.
\end{align*}

\noindent\underline{$(a,s)=(2\ell+1,1)$, $(b,t)=(2m,1)$}
\begin{align*}
  \text{LHS} &= p(2\ell+1,1)-p(2m,1)=4\ell+2-(4m-2k-2)=4\ell-4m+2k+4,\\
  \text{RHS} &= p(2\ell+1+2m-(4\ell+2),0)
  = p(-2\ell+2m-1,0)=4\ell-4m+4+2k.
\end{align*}

\noindent\underline{$(a,s)=(2\ell+1,1)$, $(b,t)=(2m+1,1)$}
\begin{align*}
  \text{LHS} &= p(2\ell+1,1)-p(2m+1,1)=4\ell+2-(4m+2)=4\ell-4m,\\
  \text{RHS} &= p(2\ell+1+2m+1-(4\ell+2),0)
  = p(-2\ell+2m,0)=-4\ell+4m.
%  \text{LHS of ($\star\star\star$)} &= 2\ell-1+2m-1\\
%  \text{RHS of ($\star\star\star$)} &= -2\ell+2m
\end{align*}
}

%%%%%%%%%%%%%%%%%%%%%%%%%%%%%%%%% 
\section*{Appendix C}
This is a suppliment for Proposition \ref{prop11} (2).
\medskip

{\tiny
\noindent\underline{$(a,s)=(2\ell,0)$, $(b,t)=(2m,0)$}
\begin{align*}
  \text{LHS} &= p(2\ell,0)+p(2m,0)=-4\ell-4m,\\
  \text{RHS} &= p(2\ell+2m,0)=-4\ell-4m.
\end{align*}

\noindent\underline{$(a,s)=(2\ell,0)$, $(b,t)=(2m+1,0)$}
\begin{align*}
  \text{LHS} &= p(2\ell,0)+p(2m+1,0)=-4\ell-4m+2k,\\
  \text{RHS} &= p(2\ell+2m+1,0)=-4\ell-4m+2k.
\end{align*}

\noindent\underline{$(a,s)=(2\ell,0)$, $(b,t)=(2m,1)$}
\begin{align*}
  \text{LHS} &= p(2\ell,0)+p(2m,1)=-4\ell+4m,\\
  \text{RHS} &= p(2\ell+2m+p(2\ell,0),1)=p(2\ell+2m-4\ell,1)=-4\ell+4m.
\end{align*}

\noindent\underline{$(a,s)=(2\ell,0)$, $(b,t)=(2m+1,1)$}
\begin{align*}
  \text{LHS} &= p(2\ell,0)+p(2m+1,1)=-4\ell+4m-2k,\\
  \text{RHS} &= p(-2\ell+2m+1,1)=-4\ell+4m-2k.
\end{align*}

\noindent\underline{$(a,s)=(2\ell+1,0)$, $(b,t)=(2m,0)$}
\begin{align*}
  \text{LHS} &= p(2\ell+1,0)-p(2m,0)=-4\ell+2k+4m,\\
  \text{RHS} &= p(2\ell+1-2m,0)=-4\ell+4m+2k.
\end{align*}

\noindent\underline{$(a,s)=(2\ell+1,0)$, $(b,t)=(2m+1,0)$}
\begin{align*}
  \text{LHS} &= p(2\ell+1,0)-p(2m+1,0)=-4\ell+2k+4m-2k=-4\ell+4m,\\
  \text{RHS} &= p(2\ell-(2m+1),0)=-4\ell+4m.
\end{align*}

\noindent\underline{$(a,s)=(2\ell+1,0)$, $(b,t)=(2m,1)$}
\begin{align*}
  \text{LHS} &= p(2\ell+1,0)-p(2m,1)=-4\ell+2k-4m,\\
  \text{RHS} &= p(2\ell+1-2m+p(2\ell+1,0),1)
  =p(-2\ell-2m+1+2k,1)=-4\ell-4m+2k.\\
\end{align*}

\noindent\underline{$(a,s)=(2\ell+1,0)$, $(b,t)=(2m+1,1)$}
\begin{align*}
  \text{LHS} &= p(2\ell+1,0)-p(2m-1,1)=-4\ell+2k-4m+2k=-4\ell-4m+4k,\\
  \text{RHS} &= p(2\ell+1-(2m+1)+p(2\ell+1,0),1)=-4\ell-4m+4k.
\end{align*}

\noindent\underline{$(a,s)=(2\ell,1)$, $(b,t)=(2m,0)$}
\begin{align*}
  \text{LHS} &= p(2\ell,1)-p(2m,0)=4\ell+4m,\\
  \text{RHS} &= p(2\ell+2m,1)=4\ell+4m.
\end{align*}

\noindent\underline{$(a,s)=(2\ell,1)$, $(b,t)=(2m+1,0)$}
\begin{align*}
  \text{LHS} &= p(2\ell,1)-p(2m+1,0)=4\ell-(-4m+p(1,0))=4\ell+4m-2k,\\
  \text{RHS} &= p(2\ell+2m+1,1)=4\ell+4m-2k. 
\end{align*}

\noindent\underline{$(a,s)=(2\ell,1)$, $(b,t)=(2m,1)$}
\begin{align*}
  \text{LHS} &= p(2\ell,1)-p(2m,1)=4\ell-4m,\\
  \text{RHS} &= p(2\ell+2m-p(2\ell,1),0)=p(-2\ell+2m)=4\ell-4m.
\end{align*}

\noindent\underline{$(a,s)=(2\ell,1)$, $(b,t)=(2m+1,1)$}
\begin{align*}
  \text{LHS} &= p(2\ell,1)-p(2m+1,1)=4\ell-4m+2k, \\
  \text{RHS} &= p(2\ell-(2m+1)-p(2\ell,1),0)
                                   =p(-2\ell+2m+1,0)=4\ell-4m+2k.
\end{align*}

\noindent\underline{$(a,s)=(2\ell+1,1)$, $(b,t)=(2m,0)$}
\begin{align*}
  \text{LHS} &= p(2\ell+1,1)+p(2m,0)=4\ell-2k-4m,\\
  \text{RHS} &= p(2\ell+1-2m,1)=4\ell-4m-2k.
\end{align*}

\noindent\underline{$(a,s)=(2\ell+1,1)$, $(b,t)=(2m+1,0)$}
\begin{align*}
  \text{LHS} &= p(2\ell+1,1)+p(2m+1,0)=4\ell-2k-4m+2k=4\ell-4m,\\
  \text{RHS} &= p(2\ell+1-(2m+1),1)=p(2\ell-2m,1)=4\ell-4m.
\end{align*}

\noindent\underline{$(a,s)=(2\ell+1,1)$, $(b,t)=(2m,1)$}
\begin{align*}
  \text{LHS} &= p(2\ell+1,1)+p(2m,1)=4\ell-2k+4m,\\
  \text{RHS} &= p(2\ell-2m-p(2\ell+1,1),0)
                                   =p(-2\ell-2m+1+2k,0)=4\ell+4m-2k.
\end{align*}

\noindent\underline{$(a,s)=(2\ell-1,1)$, $(b,t)=(2m+1,1)$}
\begin{align*}
  \text{LHS} &=  p(2\ell+1,1)+p(2m+1,1)=4\ell-2k+4m-2k=4\ell+4m-4k,\\
  \text{RHS} &= p(2\ell+1-(2m+1)-(4\ell-2k),0)
                                   =p(-2\ell-2m+2k,0)=4\ell-4m-4k.
\end{align*}
}

\end{document}
